\section{Introduction}
\label{sec:introduction}

\noindent
This paper is organized in two parts: chapter~\ref{chap:1}~goes over the preliminaries, namely, a brief introduction to topology, differential geometry, and Riemannian geometry, while chapter~\ref{chap:2}~delves into geometric inverse problems relating to medical imaging techniques.

The main purpose of this paper is to explore the mathematical foundations of the boundary rigidity problem, which is a fundamental question in Riemannian geometry with significant applications in medical imaging. The boundary rigidity problem asks whether a Riemannian manifold can be uniquely determined by the distances between points on its boundary. In the context of medical imaging, this problem is closely related to techniques such as ultrasound tomography, where one seeks to reconstruct the internal structure of a body from measurements taken at its surface. By understanding the mathematical principles underlying the boundary rigidity problem, we can gain insights into the limitations and capabilities of these imaging techniques, and potentially develop new methods for improving image reconstruction. This paper will provide a comprehensive overview of the mathematical tools and concepts necessary to approach the boundary rigidity problem, as well as a detailed analysis of its implications for medical imaging.
